% ============================================================
% CHƯƠNG 2: CƠ SỞ LÝ THUYẾT
% ============================================================
\clearpage
\section{Cơ sở lý thuyết}

\subsection{Data Warehouse và OLAP}

\subsubsection{Khái niệm Data Warehouse và Star Schema}

\textbf{Data Warehouse (Kho dữ liệu)} là hệ thống lưu trữ dữ liệu tích hợp, hướng chủ đề (subject-oriented),
không biến động (non-volatile) và biến thiên theo thời gian (time-variant), được thiết kế phục vụ cho mục đích
phân tích và hỗ trợ ra quyết định \cite{inmon2005building}.
Khác với cơ sở dữ liệu tác nghiệp (OLTP) tối ưu cho giao dịch nhanh,
Data Warehouse được tối ưu cho các truy vấn phân tích phức tạp trên khối lượng dữ liệu lớn.

\textbf{Star Schema} là mô hình tổ chức dữ liệu phổ biến nhất trong Data Warehouse,
được đề xuất bởi Ralph Kimball \cite{kimball2013datawarehouse}.
Cấu trúc bao gồm:
\begin{itemize}
    \item \textbf{Fact Table (Bảng sự kiện):} Lưu trữ các chỉ số đo lường định lượng (measures)
          của sự kiện kinh doanh, ví dụ: thời gian delay, khoảng cách bay, số lượng chuyến bay.
          Fact table chứa các khóa ngoại (foreign key) trỏ đến các bảng chiều.
    \item \textbf{Dimension Tables (Bảng chiều):} Mô tả các thuộc tính ngữ cảnh của sự kiện
          (hãng bay, sân bay, thời gian, địa lý). Mỗi bảng chiều liên kết trực tiếp với fact table,
          tạo hình dạng ngôi sao đặc trưng.
\end{itemize}

So với \textbf{Snowflake Schema} --- biến thể chuẩn hóa của Star Schema trong đó các bảng chiều
được phân rã thêm --- Star Schema có ưu điểm truy vấn đơn giản hơn (ít JOIN hơn),
hiệu suất cao hơn cho các OLAP query, và dễ hiểu hơn với người dùng cuối,
với đánh đổi là dư thừa dữ liệu (data redundancy) ở mức chấp nhận được.

\subsubsection{OLAP và các phép phân tích đa chiều}

\textbf{OLAP (Online Analytical Processing)} là công nghệ cho phép phân tích dữ liệu đa chiều
một cách linh hoạt và tương tác \cite{codd1993olap}.
Dữ liệu được tổ chức thành \textit{Data Cube} --- cấu trúc đa chiều cho phép xem xét dữ liệu
từ nhiều góc độ khác nhau đồng thời.

Các phép phân tích OLAP cơ bản bao gồm:
\begin{itemize}
    \item \textbf{Roll-up:} Tổng hợp dữ liệu lên mức cao hơn trong hệ thứ bậc,
          ví dụ từ tháng $\rightarrow$ quý $\rightarrow$ năm.
    \item \textbf{Drill-down:} Phân rã dữ liệu xuống mức chi tiết hơn,
          ví dụ từ hãng bay $\rightarrow$ từng chuyến bay cụ thể.
    \item \textbf{Slice:} Cắt lát Data Cube theo một giá trị cố định của một chiều,
          ví dụ chỉ xem dữ liệu tháng 12.
    \item \textbf{Dice:} Lọc Data Cube theo nhiều điều kiện đồng thời trên nhiều chiều.
    \item \textbf{Pivot (Rotate):} Xoay trục nhìn của Data Cube để quan sát từ góc độ khác.
\end{itemize}

% ------------------------------------------------------------
\subsection{Tiền xử lý dữ liệu}

Tiền xử lý dữ liệu là bước quan trọng nhằm đảm bảo chất lượng dữ liệu đầu vào cho cả phân tích
OLAP lẫn huấn luyện mô hình khai phá. Quy trình bao gồm bốn giai đoạn chính \cite{han2011datamining}:

\subsubsection{Data Cleaning (Làm sạch dữ liệu)}

Làm sạch dữ liệu xử lý ba vấn đề chính:
\begin{itemize}
    \item \textbf{Missing values (Giá trị thiếu):} Dữ liệu hàng không có thể thiếu thông tin
          về thời gian delay khi chuyến bay bị hủy (CANCELLED = 1) hoặc chuyển hướng (DIVERTED = 1).
          Chiến lược xử lý: điền giá trị 0 cho các cột delay khi chuyến bay bị hủy
          (không có delay thực tế), loại bỏ bản ghi không xác định được nguyên nhân.
    \item \textbf{Outliers (Giá trị ngoại lệ):} Các bản ghi có thời gian delay âm
          (tàu bay đến sớm hơn dự kiến) hoặc cực lớn (>1000 phút) cần được kiểm tra và xử lý phù hợp.
          Giá trị delay âm được giữ lại (tàu đến sớm) nhưng được clip ở ngưỡng hợp lý khi feature engineering.
    \item \textbf{Inconsistency (Mâu thuẫn):} Kiểm tra tính nhất quán giữa các cột,
          ví dụ tổng thời gian delay các thành phần phải xấp xỉ tổng delay thực tế (\texttt{ARRIVAL\_DELAY}).
\end{itemize}

\subsubsection{Data Integration (Tích hợp dữ liệu)}

Tích hợp dữ liệu từ 3 nguồn CSV độc lập:
\begin{itemize}
    \item \texttt{flights.csv} $\leftarrow$ nguồn chính chứa sự kiện chuyến bay.
    \item \texttt{airlines.csv} $\leftarrow$ tra cứu tên đầy đủ hãng bay từ mã IATA.
    \item \texttt{airports.csv} $\leftarrow$ tra cứu thông tin sân bay (tên, thành phố, tiểu bang, tọa độ).
\end{itemize}
Quá trình tích hợp sử dụng lookup join trong quá trình ETL để điền thông tin
vào các bảng chiều tương ứng, tránh lưu trùng lặp trong fact table.

\subsubsection{Data Transformation (Biến đổi dữ liệu)}

\begin{itemize}
    \item \textbf{Normalization (Chuẩn hóa):} Các cột số học như \texttt{DEPARTURE\_DELAY},
          \texttt{DISTANCE}, \texttt{SCHEDULED\_DEPARTURE} được chuẩn hóa về khoảng [0, 1]
          hoặc chuẩn hóa z-score trước khi đưa vào mô hình học máy.
    \item \textbf{Feature construction (Xây dựng đặc trưng mới):}
          Tạo các đặc trưng dẫn xuất như: giờ khởi hành (0--23), tháng trong năm, ngày trong tuần,
          nhãn phân loại \texttt{is\_delayed} (1 nếu \texttt{ARRIVAL\_DELAY} $\geq$ 15 phút, 0 nếu ngược lại).
    \item \textbf{Encoding:} Các biến phân loại (mã hãng bay, mã sân bay) được mã hóa
          bằng Label Encoding hoặc Target Encoding để phù hợp với đầu vào của các thuật toán.
\end{itemize}

\subsubsection{Data Reduction (Giảm chiều dữ liệu)}

\begin{itemize}
    \item \textbf{Attribute selection (Chọn lọc thuộc tính):} Từ 31 cột trong \texttt{flights.csv},
          chọn ra 10 đặc trưng có thể biết \textit{trước khi} chuyến bay khởi hành để đảm bảo
          mô hình có thể được sử dụng thực tế (không dùng thông tin rò rỉ tương lai).
    \item \textbf{Sampling (Lấy mẫu):} Do mất cân bằng lớp (79\% không delay, 21\% delay),
          áp dụng chiến lược class\_weight để cân bằng ảnh hưởng trong quá trình huấn luyện,
          thay vì oversampling/undersampling vật lý để tránh tăng thời gian huấn luyện.
\end{itemize}

% ------------------------------------------------------------
\subsection{Các thuật toán khai phá dữ liệu}

\subsubsection{Logistic Regression}

Logistic Regression là thuật toán phân loại tuyến tính cơ bản,
mô hình hóa xác suất thuộc lớp dương theo hàm sigmoid:
\[
P(y=1 \mid \mathbf{x}) = \frac{1}{1 + e^{-(\mathbf{w}^T\mathbf{x} + b)}}
\]
Tham số $\mathbf{w}$ được tối ưu bằng Maximum Likelihood Estimation (MLE) với regularization L1/L2.
Ưu điểm: nhanh, giải thích được qua hệ số $\mathbf{w}$, xác suất đầu ra calibrate tốt.
Hạn chế: chỉ mô hình hóa được quan hệ tuyến tính.

\subsubsection{Random Forest}

Random Forest là phương pháp ensemble kết hợp nhiều Decision Tree \cite{breiman2001random}.
Mỗi cây được huấn luyện trên bootstrap sample và chỉ xét ngẫu nhiên $\sqrt{d}$ đặc trưng tại mỗi node,
làm giảm tương quan giữa các cây. Kết quả phân loại là majority vote của tất cả cây.
Random Forest có khả năng bắt được quan hệ phi tuyến và cung cấp Feature Importance tự nhiên.

\subsubsection{Neural Network (Mạng nơ-ron nhân tạo)}

Mô hình Multilayer Perceptron (MLP) với các lớp fully-connected, hàm kích hoạt ReLU ở các lớp ẩn
và sigmoid ở lớp đầu ra cho bài toán phân loại nhị phân.
Mô hình được tối ưu bằng Adam optimizer với binary cross-entropy loss.
Neural Network có khả năng học các biểu diễn đặc trưng phức tạp nhưng đòi hỏi nhiều dữ liệu
và thời gian huấn luyện hơn.

\subsubsection{Gradient Boosting: XGBoost và LightGBM}

\textbf{XGBoost} \cite{chen2016xgboost} và \textbf{LightGBM} \cite{ke2017lightgbm} là hai triển khai
hiệu quả của thuật toán Gradient Boosting Decision Trees (GBDT).
Thuật toán xây dựng cây quyết định tuần tự, mỗi cây mới học để sửa lỗi của các cây trước:
\[
F_m(\mathbf{x}) = F_{m-1}(\mathbf{x}) + \eta \cdot h_m(\mathbf{x})
\]
trong đó $h_m$ là cây mới được tối ưu để giảm residual, $\eta$ là learning rate.
LightGBM sử dụng Leaf-wise tree growth và Gradient-based One-Side Sampling (GOSS)
giúp huấn luyện nhanh hơn XGBoost trên dữ liệu lớn.

\subsubsection{Các chỉ số đánh giá}

Với bài toán phân loại nhị phân có mất cân bằng lớp, các chỉ số đánh giá bao gồm:

\begin{itemize}
    \item \textbf{AUC-ROC:} Diện tích dưới đường ROC (Receiver Operating Characteristic),
          đo khả năng phân biệt hai lớp bất kể ngưỡng quyết định. Giá trị từ 0,5 (ngẫu nhiên) đến 1 (hoàn hảo).
    \item \textbf{F1-score:} Trung bình điều hòa của Precision và Recall:
          $F_1 = 2 \cdot \frac{\text{Precision} \times \text{Recall}}{\text{Precision} + \text{Recall}}$
    \item \textbf{Precision:} Tỷ lệ dự đoán delay đúng trong tổng số dự đoán delay.
          $\text{Precision} = \frac{TP}{TP + FP}$
    \item \textbf{Recall:} Tỷ lệ chuyến bay delay thực sự được phát hiện.
          $\text{Recall} = \frac{TP}{TP + FN}$
\end{itemize}

Trong ngữ cảnh này, Recall quan trọng hơn Precision vì bỏ sót chuyến bay delay thực sự
(False Negative) gây hại lớn hơn cảnh báo nhầm (False Positive).
